\section*{Chapter \ref{ch:s1:short-term-futures-strategies} --- Short-Term Futures Strategies}

\begin{solution}{exo:s1:short-term-futures-strategies:1}
$0.01/100 = 1$ basis point. A round trip at 1.08 ticks of spread pays half the spread twice: 1.08 basis points.
\end{solution}

\begin{solution}{exo:s1:short-term-futures-strategies:2}
$50 \times 0.3 = 15$ basis points. Eight events a year multiply a per-event Sharpe ratio by $\sqrt{8}$; with the chapter's spacing of 32 trading days,
$\sqrt{252/32} = 2.81$.
\end{solution}

\begin{solution}{exo:s1:short-term-futures-strategies:3}
The planted trend drift is spread evenly over the session, so the longer the range, the more of the day's direction it has seen when the breakout
fires, and the breakout trades at most once a day, so the later entry costs little.
\end{solution}

\begin{solution}{exo:s1:short-term-futures-strategies:4}
$1\%/\sqrt{39} = 16$ basis points; a true 50 basis points gives $t = 0.5/0.16 = 3.1$ on average (2.7 in the sample).
\end{solution}

\begin{solution}{exo:s1:short-term-futures-strategies:5}
Its gross edge per trade, the part of a move that bounces back, is smaller than half a spread in and half a spread out; with several trades a day the
costs grow with the number of trades, the edge does not.
\end{solution}

\begin{solution}{exo:s1:short-term-futures-strategies:6}
The session's daily sum is scaled by about $(1 - 0.1) = 0.9$, so its variance by 0.81; with three quarters of the variance in the session, the day's
volatility is multiplied by $\sqrt{0.25 + 0.75 \times 0.81} = 0.926$.
\end{solution}

\begin{solution}{exo:s1:short-term-futures-strategies:7}
At a threshold of two standard deviations the fade earns 5.18 basis points a day before costs on 7.3 trades a day: it breaks even at 0.35 basis points
per side, below the tape's 0.54. A passive entry at the far side of the spread on one leg would bring it near break-even; both legs passive, with fills
uncertain, is what a liquidity provider does.
\end{solution}

\begin{solution}{exo:s1:short-term-futures-strategies:8}
Forty events have a standard error of about 16 basis points: a mean of 40 is consistent with a true drift anywhere from about 10 to 70, including one
decayed to a third. Five years cannot confirm that the edge is intact; the paper's prior, the costs and how crowded the trade has become matter more.
\end{solution}

\begin{solution}{pb:s1:short-term-futures-strategies:1}
\begin{enumerate}
\item One-minute bars over ten years: U-shaped volatility, overnight gap, daily trend drift, bounce, and a \omterm{def:s1:short-term-futures-strategies:pre}{pre-announcement drift} cut after year five.
\item Half the time-weighted quoted spread of a tape session: 0.54 basis points.
\item A position on the first move outside the first minutes' range; fading short moves that bounce back.
\item The trend day; the bounce.
\item 0.35, 0.55 and 1.25 after costs for 15, 30 and 60 minutes.
\item 3.78, 2.86, 1.82 and 0.59 before; $-3.31$, $-1.50$, $-0.71$ and $-0.78$ after, for thresholds of 1.5 to 3.
\item The breakout trades once a day for a large edge; the fade trades often for a small one.
\item The liquidity providers who earn the spread.
\item A move before a scheduled announcement whatever its content; large US equity returns before FOMC decisions, none in Treasuries or before other
announcements.
\item The eight 2026 FOMC meetings.
\item 40.6 and 18.8 basis points per event, Sharpe ratios of 1.21 and 0.58.
\item Eight events a year give a standard error of about 16 basis points in five years.
\item \textbf{Named result.} The pre-announcement hold earned 40.6 basis points per event after costs ($t = 2.7$) and a Sharpe ratio of 1.21 before the
planted crowding, and 18.8 basis points ($t = 1.3$) and 0.58 after it.
\item Turnover against edge, sample size against decay, and the mechanism.
\item Scheduled meeting dates only, statement times as scheduled, costs, and enough years to measure decay.
\item Post-announcement momentum.
\item Chapter~13 measured where the day's return and volume fall; this chapter trades inside the session on those patterns.
\item Passive fills and a lower cost than the bounce's size, which means being a liquidity provider.
\item It has a sign whatever the decision; information would move prices in the direction of the news.
\item Immediacy and risk around predictable moments.
\end{enumerate}
\end{solution}

\begin{solution}{iq:s1:short-term-futures-strategies:1}
It takes a position in the direction of the first move outside the session's first minutes' range; it fails on range days, when breakouts reverse,
and when many traders' stops sit at the same levels.
\end{solution}

\begin{solution}{iq:s1:short-term-futures-strategies:2}
Measure it only after the publication date, with the same definition; compare its size with the published one and its standard error; check other
markets and other announcements as controls; and check whether positioning before meetings has grown.
\end{solution}

\begin{solution}{iq:s1:short-term-futures-strategies:3}
Costs: the number of trades times the spread and fees against the gross edge per trade; then whether fills are realistic (queue position, adverse
selection) and whether prices used were tradeable.
\end{solution}

\begin{solution}{iq:s1:short-term-futures-strategies:4}
Bars or messages with exact timestamps, the spread at each trade, a fill model for passive orders, session calendars and rolls, scheduled \omterm{def:s1:earnings:event}{event times},
and latency between signal and order.
\end{solution}

\begin{solution}{iq:s1:short-term-futures-strategies:5}
A jump at the announcement against the position, wider spreads and thinner books around it, and a crowded exit if many hold the same position.
\end{solution}

\begin{solution}{iq:s1:short-term-futures-strategies:6}
Holding $-\operatorname{sign}(r_t)$ over the next minute earns $-\operatorname{sign}(r_t)(e_{t+1} + \theta e_t)$, with expectation
\[
-\theta\,E[\operatorname{sign}(r_t)\,e_t] = -\theta\sigma\sqrt{2/\pi}\big/\sqrt{1 + \theta^2}
\]
for Gaussian shocks, positive for $\theta < 0$. With a cost $k$ each way it breaks even when $2k$ equals that expectation: for $\theta = -0.1$ the edge
is about 8\% of a minute's volatility, which must exceed a spread.
\end{solution}
